\subsection{Environment overview}

We study a repeated public-goods economy with heterogeneous country roles.
Each agent is labeled developed or developing.
Labels govern initial wealth, vulnerability, and institutional assignment in the climate package.
Interaction proceeds in discrete rounds.
Within each round, agents choose Stage-1 contributions inside their assigned institution.
Agents in the enforcing institution then allocate Stage-2 sanctions and rewards.
Climatic shocks, fund accounting, social scoring, and democracy occur on a fixed schedule described in Section~4.

Two institutions operate in parallel.
The enforcing institution (SI) combines Stage-1 public-goods returns with peer punishment and reward.
The non-enforcing institution (SFI) provides Stage-1 returns without Stage-2 sanctioning.
Institution labels are presented to agents in climate language as a binding treaty versus a non-binding agreement.
In the primary climate run, assignment is forced every round.
Developed agents enter SI.
Developing agents enter SFI.
Free institutional choice is reserved for boundary baselines discussed later.

\subsection{Stage-1 public good}

Let $c_i$ denote agent $i$'s Stage-1 contribution and let $G$ denote the set of agents in the same institution.
Total contribution is $C=\sum_{j\in G}c_j$.
With public-good multiplier $m$ and group size $n=|G|$, each member receives Stage-1 share
\begin{equation}
\mathrm{share}_i=\frac{m\cdot C}{n}.
\end{equation}
The marginal per-capita return is therefore $m/n$.
Contribution budgets equal current wealth in the climate package, floored at zero.
Agents receive decision cards that state budget bounds, group size, marginal return, and recent same-institution averages.

Contribution intensity is the primary behavioral outcome.
For each agent-round we define
\begin{equation}
\mathrm{prop}_{i,t}=\frac{c_{i,t}}{w_{i,t}},
\end{equation}
where $w_{i,t}$ is end-of-round wealth.
Absolute contributions differ by orders of magnitude across groups.
Intensity normalizes for scale and supports cross-agent event studies.

\subsection{Stage-2 enforcement}

Only SI members enter Stage-2.
Each receives a sanction budget proportional to wealth, with a climate floor.
Punishment tokens reduce a target's payoff at a fixed effect-to-cost ratio.
Reward tokens raise a target's payoff at unit cost.
Targets are listed with current contributions and pairwise consistency scores when available.
Unspent Stage-2 budget is retained according to the institutional rules presented in the decision card.
SFI members skip Stage-2 entirely.

Enforcement is itself a second-order public good.
High contributors may bear disproportionate sanction spending.
We track that burden as context.
We do not treat it as the primary estimand.

\subsection{Dual-use loss-and-damage fund}

Stage-1 contributions are dual-use deposits.
Each contributed unit also enters a collective loss-and-damage pool.
The pool is hidden from agents as a live balance.
Agents are told that contributions support both emissions-reduction framing and disaster payouts.
After climatic shocks, developing agents may receive payouts subject to coverage caps and equity weights.
Developed agents absorb damage without fund receipts under the default disbursement rule.

Gross damage for agent $i$ after shock severity $s$ is proportional to $s$ and vulnerability $v_i$.
Net damage equals gross damage minus any fund receipt.
Coverage is the ratio of aggregate payouts to aggregate gross damage in a shock round.
High coverage can coexist with widening wealth gaps when collection and growth dominate disbursement.

\subsection{Social information channels}

Three linked channels create social pressure.

\paragraph{Pairwise consistency scoring.}
After contributions, evaluators score targets on a one-to-ten scale.
High scores indicate match between stated intent and action.
Low scores indicate inconsistency.
Scores are elicited with short natural-language justifications.

\paragraph{Aggregate reputation.}
Agents observe an aggregate peer-trust rating derived from pairwise scores.
Bad reputation is defined analytically as an aggregate score below four.
A reputation drop is a decline of at least one point between rounds.

\paragraph{Gossip bulletins.}
Low pairwise scores generate gossip items injected into subsequent prompts.
Bulletins name the observer, the target, the consistency score, and a short comment.
Capacity limits keep only a top set of items.
In the event study, gossip targeting is reconstructed from low-score conditions consistent with bulletin generation.
The simulation does not store a separate verbatim gossip transcript for every edge.

These channels are informational.
They do not mechanically force higher contributions.
Any repair must arise from prompted choice.

\subsection{Democracy}

Every five rounds, agents may propose and vote on a bounded whitelist of parameters.
Whitelist items include subsidy fraction, subsidy breadth, punishment and reward effects, and loss-and-damage equity or damage weights.
Majority outcomes update the live rule set.
Democracy is endogenous institutional adaptation.
It is analyzed here as backdrop for contribution dynamics.

\subsection{Agent decision interface}

Agents are prompted language models with a default self-interested persona.
They maximize cumulative payoff unless a specialized persona is assigned.
Climate role guidance soft-frames developed versus developing constraints.
Decision outputs are structured as contribution amounts, sanction maps, proposals, votes, and rationales.
Belief states update after rounds with trust labels and institutional strategies.
We analyze both numeric choices and coded language.
We never equate fluent rationales with internalized norms.
